\documentclass[11pt]{beamer}

% Theme
\usetheme{Madrid}
\usecolortheme{beaver}
\usefonttheme{professionalfonts}

% Packages
\usepackage[T1]{fontenc}
\usepackage[utf8]{inputenc}
\usepackage{lmodern}
\usepackage{amsmath,amssymb,amsthm}
\usepackage{graphicx}
\usepackage{hyperref}
\usepackage{tikz}
\usetikzlibrary{arrows,positioning}
\usepackage{listings}
\usepackage{caption}
\usepackage{subcaption}
\usepackage{epstopdf}
\usepackage{mathrsfs}
\usepackage{booktabs}
\usepackage{pgfplots}
\pgfplotsset{compat=1.17}

% Bibliography (uncomment refs frame at bottom to show)
\usepackage{biblatex}
\addbibresource{refs.bib}

% Code style
\lstset{
  basicstyle=\ttfamily\small,
  keywordstyle=\color{blue},
  commentstyle=\color{gray},
  stringstyle=\color{red!60!black},
  showstringspaces=false,
  breaklines=true,
  frame=single
}

% Title
%\title[Lecture No.]{Topic Name}
%\author[name1\\name2]{\textcolor{blue}{Name1 \\ Name 2}\\ \scriptsize Dept. of CSE, IIT Bombay}
%\institute[IIT Bombay]{IIT Bombay}
% \date{\small July 31, 2025}



% Title
\title[CS6007 Lecture: Lecture Name]{Presentation on Topic Name}

% Title
\title[CS6007 Lecture]{Presentation on Topic Name}

% Show ONLY names in the bottom-left footer via the optional [] argument
\author[Name 1, Name 2]{%
  \textbf{Presented by:}\\
  \begin{tabular}{rl}
    \textcolor{blue}{Abhay Chandra(roll no.)} & Dept. of CSE \\
    \textcolor{blue}{Chaitanya De(roll no.)} & Dept. of CSE \\
  \end{tabular}\\[1em]
  \textbf{Course:} CS6007:Multi-Agent Machine Learning\\[0.3em]
  \textbf{Instructor:} Prof.\ Avishek Ghosh
}

\institute[IIT Bombay]{IIT Bombay}
% \date{\small August 2025}




% Small logo on each frametitle (top-right)
\addtobeamertemplate{frametitle}{}{%
  \begin{tikzpicture}[remember picture,overlay]
    \node[anchor=north east, xshift=0.2cm, yshift=0.4cm] at (current page.north east) {%
      \includegraphics[height=1.2cm]{iitb_logo.png}%
    };
  \end{tikzpicture}%
}

\begin{document}

% Title slide with faint logo background
\usebackgroundtemplate{%
  \tikz[remember picture,overlay] \node[opacity=0.2] at ([yshift=-2cm]current page.center) {%
    \includegraphics[width=0.6\paperwidth]{iitb_logo.png}%
  };
}
\begin{frame}
  \titlepage
\end{frame}
\usebackgroundtemplate{} % clear background

% Outline
\begin{frame}{Outline}
  \setbeamertemplate{section in toc}[sections numbered]
  \tableofcontents
\end{frame}

%--------------------------------------------------
% Content

\section{Introduction}

\begin{frame}{Motivation}
  \begin{itemize}
    \item Content recommendation is crucial for online services
    \item Users have diverse preferences and interact with different content
    \item Challenge: Identify clusters of similar users with unknown group structure
    \item Need adaptive algorithms that scale to large numbers of users
  \end{itemize}
\end{frame}

\begin{frame}{Problem Statement}
  \begin{itemize}
    \item Partition $n$ users into $m$ clusters: $\mathcal{V} = V_1, V_2, \ldots, V_m$ (unknown)
    \item Users within each cluster share similar preferences
    \item At each round $t$, a user is served with context vectors
    \item Goal: Design algorithm to learn clustering and optimize recommendations
    \item Trade-off between exploration and exploitation
  \end{itemize}
\end{frame}

\begin{frame}{Key Challenges}
  \begin{itemize}
    \item \textbf{Implicit network:} Social structure is encoded in unknown clustering
    \item \textbf{Adaptive clustering:} Users' preferences may shift over time
    \item \textbf{Limited feedback:} Only payoff signals from served users are observed
    \item \textbf{Scalability:} Algorithm must handle large numbers of users and features
    \item \textbf{Regret bounds:} Achieve theoretical guarantees with high probability
  \end{itemize}
\end{frame}

\section{Learning Model}

\begin{frame}{Problem Formulation}
  \begin{itemize}
    \item Users: $V = \{1, \ldots, n\}$ partitioned into $m$ clusters $V_1, \ldots, V_m$
    \item Each cluster $j$ has unknown parameter vector $u_j \in \mathbb{R}^d$
    \item At round $t$, user $i_t$ is served with context vectors $C_{i_t} = \{x_1, \ldots, x_s\}$ from $\mathbb{R}^d$
    \item Payoff: $o_t(x) = u_{j(i_t)} \cdot x + \eta_{j(i_t)}(x)$ where $\eta$ is bounded noise
    \item Assume well-separated clusters: $\|u_i - u_j\| > \gamma > 0$ for $i \neq j$
  \end{itemize}
\end{frame}

\begin{frame}{Model Assumptions}
  \begin{enumerate}
    \item User behavior encodes unknown clustering with users in same cluster having similar preferences
    \item Context vectors in cluster $C_{i_t}$ generated from ball surface near mean
    \item Payoff assumptions: linear in cluster vectors with bounded additive noise
    \item Well-separated clusters ensure identifiability
    \item Sequential learning: At each round, algorithm observes user, context, and selected payoff
  \end{enumerate}
\end{frame}

\section{Algorithm: CLUB}

\begin{frame}{CLUB Algorithm Overview}
  \textbf{Cluster of Bandits (CLUB)} - Key Ideas:
  \begin{itemize}
    \item Maintain estimate $w_{i,t}$ for each user's preference vector
    \item Detect potential cluster merges using \textbf{confidence bounds}
    \item Use \textbf{least-squares} estimation for cluster vectors
    \item Adaptively update cluster structure based on observed similarity
    \item Efficient update: $O(d^2)$ per round
  \end{itemize}
\end{frame}

\begin{frame}{CLUB Algorithm: Pseudocode (Part 1)}
  \textbf{Initialization:}
  \begin{itemize}
    \item Initialize estimate vectors $w_{i,0} = 0 \in \mathbb{R}^d$, matrices $M_{i,0} = I$
    \item Clusters initially: $\mathcal{V}_{1,1} = V$, number of clusters $m_1 = 1$
  \end{itemize}
  
  \textbf{For each round $t = 1, 2, \ldots, T$:}
  \begin{enumerate}
    \item Receive user $i_t \in V$ and context vectors $C_t = \{x_1, \ldots, x_s\}$
    \item Select context $x_t$ to recommend (bandit selection)
    \item Observe payoff $o_t \in [-1, 1]$
    \item Update estimates and confidence bounds
  \end{enumerate}
\end{frame}

\begin{frame}{CLUB Algorithm: Pseudocode (Part 2)}
  \textbf{Update Phase:}
  \begin{itemize}
    \item Compute confidence bound $\text{CB}_{i,t-1}(x)$ for each user
    \item Update matrix: $M_{i,t} = M_{i,t-1} + x x^T$
    \item Update vector: $b_{i,t} = b_{i,t-1} + o_t x$
    \item Least-squares estimate: $w_{i,t} = M_{i,t}^{-1} b_{i,t}$
  \end{itemize}
  
  \textbf{Clustering Step:}
  \begin{itemize}
    \item For each pair of nodes, check if $\|w_{i,t-1} - w_{i',t-1}\|$ exceeds threshold
    \item If distance is small: merge clusters
    \item Update aggregate weight vectors for clusters
    \item Maintain graph of connected components
  \end{itemize}
\end{frame}

\begin{frame}{Key Algorithm Properties}
  \begin{itemize}
    \item \textbf{Greedy clustering:} Clusters grow from initial singletons by detecting similarity
    \item \textbf{Confidence-based merging:} Uses concentration inequalities to decide merges
    \item \textbf{Adaptive updates:} Parameter estimates improve as more data arrives
    \item \textbf{Efficient:} Maintains only necessary cluster information
    \item \textbf{Regret-aware:} Design balances learning true structure and optimizing payoffs
  \end{itemize}
\end{frame}

\section{Theoretical Analysis}

\begin{frame}{Regret Bound}
  \textbf{Main Result (Theorem 1):}
  
  With high probability, the CLUB algorithm achieves cumulative regret:
  $$\sum_{t=1}^{T} r_t = \tilde{O}\left(\left(\sigma d + \sqrt{d}\right) \sqrt{T} \left(1 + \sum_{j=1}^{m} \sqrt{\frac{|V_j|}{n}}\right)\right)$$
  
  where:
  \begin{itemize}
    \item $d$ = number of features
    \item $T$ = number of rounds
    \item $\sigma$ = noise level
    \item $|V_j|$ = size of cluster $j$
    \item $\tilde{O}$ hides poly-logarithmic factors
  \end{itemize}
\end{frame}

\begin{frame}{Regret Analysis Components}
  The regret bound consists of three main terms:
  
  \begin{enumerate}
    \item \textbf{Feature complexity term:} $(\sigma d + \sqrt{d})\sqrt{T}$
      \begin{itemize}
        \item Reflects difficulty of learning $d$-dimensional vectors
        \item Noise contributes factor of $\sigma$
      \end{itemize}
    
    \item \textbf{Transient term:} Convergence to minimal eigenvalue
      \begin{itemize}
        \item Cost of learning cluster structure
        \item Depends on geometric properties
      \end{itemize}
    
    \item \textbf{Cluster size term:} $1 + \sum_{j=1}^{m} \sqrt{|V_j|/n}$
      \begin{itemize}
        \item Small clusters incur more regret
        \item Balanced clusters are favorable
      \end{itemize}
  \end{enumerate}
\end{frame}

\begin{frame}{Key Theoretical Insights}
  \begin{itemize}
    \item \textbf{High-probability analysis:} Uses concentration inequalities and union bounds
    \item \textbf{Dimension-dependent complexity:} Regret scales with $\sqrt{d}$, not $d$
    \item \textbf{Compared to state-of-the-art:} CLUB achieves regret improvement by factor of $\sqrt{m}$
    \item \textbf{Clustering correctness:} Algorithm correctly identifies true clusters with high probability
    \item \textbf{Lemma 1:} Connectivity of random graphs ensures clusters can be detected
  \end{itemize}
\end{frame}

\section{Experiments}

\begin{frame}{Experimental Setup}
  \textbf{Datasets:}
  \begin{itemize}
    \item \textbf{Synthetic:} Controlled experiments with known clusters and varying sizes
    \item \textbf{LastFM:} Music streaming service, 1,892 users, 17,632 items
    \item \textbf{Delicious:} Social bookmarking, diverse payoff structures
    \item \textbf{Yahoo:} News recommendations, real-world scale and complexity
  \end{itemize}
  
  \textbf{Baselines:}
  \begin{itemize}
    \item LinUCB-ONE: Single model for all users
    \item LinUCB-IND: Independent models per user
    \item GOBLIN: Idealized algorithm with known clusters
    \item CLAIRVOYANT: Oracle algorithm
  \end{itemize}
\end{frame}

\begin{frame}{Results on Synthetic Data}
  \begin{itemize}
    \item CLUB consistently outperforms LinUCB-ONE and LinUCB-IND
    \item Performance gap widens as clusters become more balanced
    \item Algorithm scales to larger numbers of clusters
    \item Robustness to noise levels and feature dimensions
    \item Faster convergence with larger cluster sizes
  \end{itemize}
\end{frame}

\begin{frame}{Results on Real Datasets}
  \begin{itemize}
    \item \textbf{LastFM:} CLUB outperforms baselines on cumulative regret metric
    \item \textbf{Delicious:} Competitive performance in "many niches" scenario
    \item \textbf{Yahoo CTR:} Demonstrates clustering benefit on large-scale data
    \item Algorithm captures actual user interest similarities
    \item Parameter tuning via cross-validation improves performance
  \end{itemize}
\end{frame}

\begin{frame}{Experimental Insights}
  \begin{itemize}
    \item Clustering approach is beneficial when users naturally group by preference
    \item Performance depends on cluster structure and separation
    \item Algorithm adapts well to different data characteristics
    \item Real-world datasets show consistent improvements
    \item Trade-off between exploration and cluster learning affects regret
  \end{itemize}
\end{frame}

\section{Conclusion}

\begin{frame}{Summary}
  \textbf{Key Contributions:}
  \begin{enumerate}
    \item Novel algorithmic approach to online clustering of bandits
    \item Achieves $\tilde{O}(\sqrt{T})$ regret with clustering benefit
    \item Theoretical analysis with high-probability guarantees
    \item Empirical validation on synthetic and real datasets
    \item Practical algorithm with efficient updates
  \end{enumerate}
\end{frame}

\begin{frame}{Advantages and Limitations}
  \textbf{Advantages:}
  \begin{itemize}
    \item Provable theoretical guarantees
    \item Automatic cluster discovery
    \item Scalable to large user populations
  \end{itemize}
  
  \textbf{Limitations:}
  \begin{itemize}
    \item Assumes well-separated clusters
    \item Requires known cluster count (variant with data-dependent separation possible)
    \item Performance sensitive to feature dimensionality
  \end{itemize}
\end{frame}

\begin{frame}{Future Directions}
  \begin{itemize}
    \item Extend to non-linear payoff functions
    \item Investigate privacy-preserving clustering algorithms
    \item Apply to other recommendation domains (video, ads, etc.)
    \item Dynamic cluster updates for evolving user preferences
    \item Integration with other machine learning techniques
  \end{itemize}
\end{frame}









% External slide file (thank-you)
\input{Thank_you_Slide}

% References
% \begin{frame}[allowframebreaks]{References}
%   \printbibliography
% \end{frame}

\end{document}
